\documentclass[10pt,a4paper]{amsart}
\usepackage{amsmath,amssymb,amsthm,graphicx}
% Retained minimal preamble for the Stacks Project chapter more-algebra.tex (More on Algebra).
% The chapter file pinned-original-source.tex carries only \input{preamble} on its first line, so no
% preamble text lives inside the chapter itself.  The Stacks Project uses one single shared preamble for
% every chapter; the snapshot stages it identically as sources/stacks-chow/preamble.tex and
% sources/stacks-etale-cohomology/preamble.tex.  The block below is copied mechanically, line for line,
% from sources/stacks-chow/preamble.tex lines 14-18 (the Xy-pic setup) and 154-217 (the native theorem
% declarations together with the mathematical macros).  Nothing was added, removed or redefined; the
% wrapper class and the amsmath/amssymb/amsthm/graphicx base packages are supplied by the assembler.
% Chapter: more-algebra.tex -- More on Algebra.
% For commutative diagrams we use Xy-pic
\usepackage[all]{xy}

% We use 2cell for 2-commutative diagrams.
\xyoption{2cell}

% Theorem environments.
%
\theoremstyle{plain}
\newtheorem{theorem}[subsection]{Theorem}
\newtheorem{proposition}[subsection]{Proposition}
\newtheorem{lemma}[subsection]{Lemma}

\theoremstyle{definition}
\newtheorem{definition}[subsection]{Definition}
\newtheorem{example}[subsection]{Example}
\newtheorem{exercise}[subsection]{Exercise}
\newtheorem{situation}[subsection]{Situation}

\theoremstyle{remark}
\newtheorem{remark}[subsection]{Remark}
\newtheorem{remarks}[subsection]{Remarks}

\numberwithin{equation}{subsection}

% Macros
%
\def\lim{\mathop{\mathrm{lim}}\nolimits}
\def\colim{\mathop{\mathrm{colim}}\nolimits}
\def\Spec{\mathop{\mathrm{Spec}}}
\def\Hom{\mathop{\mathrm{Hom}}\nolimits}
\def\Ext{\mathop{\mathrm{Ext}}\nolimits}
\def\SheafHom{\mathop{\mathcal{H}\!\mathit{om}}\nolimits}
\def\SheafExt{\mathop{\mathcal{E}\!\mathit{xt}}\nolimits}
\def\Sch{\mathit{Sch}}
\def\Mor{\mathop{\mathrm{Mor}}\nolimits}
\def\Ob{\mathop{\mathrm{Ob}}\nolimits}
\def\Sh{\mathop{\mathit{Sh}}\nolimits}
\def\NL{\mathop{N\!L}\nolimits}
\def\CH{\mathop{\mathrm{CH}}\nolimits}
\def\proetale{{pro\text{-}\acute{e}tale}}
\def\etale{{\acute{e}tale}}
\def\QCoh{\mathit{QCoh}}
\def\Ker{\mathop{\mathrm{Ker}}}
\def\Im{\mathop{\mathrm{Im}}}
\def\Coker{\mathop{\mathrm{Coker}}}
\def\Coim{\mathop{\mathrm{Coim}}}

% Boxtimes
%
\DeclareMathSymbol{\boxtimes}{\mathbin}{AMSa}{"02}

%
% Macros for moduli stacks/spaces
%
\def\QCohstack{\mathcal{QC}\!\mathit{oh}}
\def\Cohstack{\mathcal{C}\!\mathit{oh}}
\def\Spacesstack{\mathcal{S}\!\mathit{paces}}
\def\Quotfunctor{\mathrm{Quot}}
\def\Hilbfunctor{\mathrm{Hilb}}
\def\Curvesstack{\mathcal{C}\!\mathit{urves}}
\def\Polarizedstack{\mathcal{P}\!\mathit{olarized}}
\def\Complexesstack{\mathcal{C}\!\mathit{omplexes}}
% \Pic is the operator that assigns to X its picard group, usage \Pic(X)
% \Picardstack_{X/B} denotes the Picard stack of X over B
% \Picardfunctor_{X/B} denotes the Picard functor of X over B
\def\Pic{\mathop{\mathrm{Pic}}\nolimits}
\def\Picardstack{\mathcal{P}\!\mathit{ic}}
\def\Picardfunctor{\mathrm{Pic}}
\def\Deformationcategory{\mathcal{D}\!\mathit{ef}}
\usepackage{cleveref}
\crefname{originallemma}{original lemma}{original lemmas}
\Crefname{originallemma}{Original lemma}{Original lemmas}
\makeatletter
\newlabel{algebra-lemma-cover}{{00EO}{source}{}{sourcepointer-1}{}}
\newlabel{algebra-lemma-cover@cref}{{[originallemma][0][]00EO}{[1][0][]source}{}{}{}}
\makeatother
\crefname{originallemma}{original lemma}{original lemmas}
\Crefname{originallemma}{Original lemma}{Original lemmas}
\makeatletter
\newlabel{algebra-lemma-finite-projective}{{00NX}{source}{}{sourcepointer-2}{}}
\newlabel{algebra-lemma-finite-projective@cref}{{[originallemma][0][]00NX}{[1][0][]source}{}{}{}}
\makeatother
\title{An invertible module is the same thing as a finite locally free module of rank one, and its inverse is its dual}
\author{The Stacks Project Authors}
\date{stacks-project:more-algebra.tex:local-snapshot-20261012; GNU FDL 1.2 or later}
\begin{document}
\maketitle
\setcounter{subsection}{1}
\noindent\textit{Editorial note.} All delivered mathematical text is the author's own, copied from the pinned Stacks Project chapter file more-algebra.tex; no statement or proof step was rewritten, reordered or removed.\par\medskip
\noindent\textit{Editorial note.} The chapter file pinned-original-source.tex carries only \textbackslash{}input\{preamble\} on its first line, so no preamble text lives inside the chapter. The retained file sources/stacks-more-algebra/retained-preamble.tex is a mechanical, line-for-line copy of the project's single shared preamble as staged in the snapshot; nothing was added, removed or redefined and no package was added.\par\medskip
\noindent\textit{Editorial note.} The source URL is the source-level Stacks Project URL, as required when more than one card is produced from one source; the per-claim tag is carried by claim\_key.\par\medskip
\noindent\textit{Editorial note.} The author writes '(details omitted)' once, at proof line 36133, for the verification that the evaluation map becomes an isomorphism after localizing; that phrase is kept verbatim and the step is immediately closed by the subsequent cover argument.\par\medskip
\noindent\textit{Allowed standard prerequisites (editorial).}\begin{itemize}
\item Finite locally free modules and the cover argument for checking an isomorphism locally. used at proof lines 36131-36134 and 36162-36163. Lets the evaluation map be verified after localizing, and characterises finite local freeness.
\item Invertible modules and tensor functors forming a quasi-inverse pair. condition (2) at statement line 36122 and used at proof lines 36138-36141. Is the equivalence with the existence of an inverse module.
\item Automorphisms of a module factoring through a finite free module. used at proof lines 36148-36161. Produces the direct summand presentation of M that yields finite local freeness.
\end{itemize}\medskip
\ifdefined\hypertarget\hypertarget{sourcepointer-1}{}\fi
\noindent\textit{Source reference.} Original lemma 00EO: focused-seeds/AGP/sources/stacks-snapshot-20261011/tags.txt, source lines 542--542. Cross-chapter navigation pointer into the Stacks Project; the standard fact it supplies is applied inside the counted proof span and the argument resting on it stays there.\par\medskip
\ifdefined\hypertarget\hypertarget{sourcepointer-2}{}\fi
\noindent\textit{Source reference.} Original lemma 00NX: focused-seeds/AGP/sources/stacks-snapshot-20261011/tags.txt, source lines 875--875. Cross-chapter navigation pointer into the Stacks Project; the standard fact it supplies is applied inside the counted proof span and the argument resting on it stays there.\par\medskip
\begin{lemma}
\label{lemma-invertible}
Let $R$ be a ring. Let $M$ be an $R$-module. Equivalent are
\begin{enumerate}
\item $M$ is finite locally free module of rank $1$,
\item $M$ is invertible, and
\item there exists an $R$-module $N$ such that $M \otimes_R N \cong R$.
\end{enumerate}
Moreover, in this case the module $N$ in (3) is isomorphic
to $\Hom_R(M, R)$.
\end{lemma}

\begin{proof}
Assume (1). Consider the module $N = \Hom_R(M, R)$ and the evaluation
map $M \otimes_R N = M \otimes_R \Hom_R(M, R) \to R$. If $f \in R$
such that $M_f \cong R_f$, then the evaluation map becomes an isomorphism
after localization at $f$ (details omitted). Thus we see the evaluation
map is an isomorphism by Algebra, Lemma \ref{algebra-lemma-cover}.
Thus (1) $\Rightarrow$ (3).

\medskip\noindent
Assume (3). Then the functor $K \mapsto K \otimes_R N$ is a quasi-inverse
to the functor $K \mapsto K \otimes_R M$. Thus (3) $\Rightarrow$ (2).
Conversely, if (2) holds, then $K \mapsto K \otimes_R M$ is essentially
surjective and we see that (3) holds.

\medskip\noindent
Assume the equivalent conditions (2) and (3) hold. Denote
$\psi : M \otimes_R N \to R$ the isomorphism from (3).
Choose an element $\xi = \sum_{i = 1, \ldots, n} x_i \otimes y_i$
such that $\psi(\xi) = 1$. Consider the isomorphisms
$$
M \to M \otimes_R M \otimes_R N \to M
$$
where the first arrow sends $x$ to $\sum x_i \otimes x \otimes y_i$
and the second arrow sends $x \otimes x' \otimes y$ to $\psi(x' \otimes y)x$.
We conclude that $x \mapsto \sum \psi(x \otimes y_i)x_i$ is
an automorphism of $M$. This automorphism factors as
$$
M \to R^{\oplus n} \to M
$$
where the first arrow is given by
$x \mapsto (\psi(x \otimes y_1), \ldots, \psi(x \otimes y_n))$
and the second arrow by $(a_1, \ldots, a_n) \mapsto \sum a_i x_i$.
In this way we conclude that $M$ is a direct summand of a finite free
$R$-module. This means that $M$ is finite locally free
(Algebra, Lemma \ref{algebra-lemma-finite-projective}).
Since the same is true for $N$ by symmetry and since
$M \otimes_R N \cong R$, we see that
$M$ and $N$ both have to have rank $1$.
\end{proof}

\end{document}
